\documentclass{article}
\usepackage{amsmath, amssymb, amsthm}

\setlength{\parindent}{0pt}

\newtheorem{claim}{Claim}

\begin{document}

\begin{claim}
    Among any $n^2 + 1$ points within an $n \times n$ square, there must exist two points whose distance 
    is at most $\sqrt{2}$.
\end{claim}

\begin{proof}
    Partition the $n \times n$ square into $n^2$ unit squares. Since there are $n^2 + 1$ points and only 
    $n^2$ unit squares, the pigeonhole principle implies that at least two of the points lie in the same 
    unit square.

    \medskip

    The greatest possible distance between two points in a $1 \times 1$ square is the length of its 
    diagonal, which is
    \[
    \sqrt{1^2+1^2}=\sqrt{2}.
    \]

    \medskip

    Hence among any $n^2 + 1$ points within an $n \times n$ square, there must exist two points whose 
    distance is at most $\sqrt{2}$.
\end{proof}

\begin{claim}
    Jellybeans of $6$ flavours are distributed among $5$ jars, with $11$ jellybeans of each flavour. 
    Then some jar contains at least three jellybeans of two distinct flavours.
\end{claim}

\begin{proof}
    Suppose for contradiction that no jar contains two distinct flavours each appearing at least three
    times. Then every jar contains at most one flavour appearing at least three times. Assign such a
    flavour to each jar whenever one exists. Since at most five flavours can be assigned in this way,
    there is a flavour that is assigned to no jar. This flavour therefore appears fewer than three times 
    in every jar, so it appears at most $2 \cdot 5 = 10$ times in total, contradicting the fact that each 
    flavour appears $11$ times.

    \medskip

    Hence some jar must contain at least three jellybeans of two distinct flavours.
\end{proof}

\begin{claim}
    Among every set of $30$ integers, there exist two whose difference or sum is a multiple of $51$.
\end{claim}

\begin{proof}
    Let $a_1, a_2, \dots a_{30}$ denote the $30$ integers and $r_i$ denote the residue of 
    $a_i \pmod{51}$, where $0 \le r_i \le 50$. Partition the residue classes modulo $51$ into the 
    following $26$ sets:
    \[
    \{0\}, \{1, 50\}, \{2, 49\}, \dots, \{25, 26\}
    \]
    Since there are $30$ residues but only $26$ sets, the pigeonhole principle guarantees that two 
    residues, say $r_x$ and $r_y$, belong to the same set.

    \medskip

    If $r_x = r_y$, then
    \[
    a_x \equiv a_y \pmod{51}
    \]
    so
    \[
    51 \mid (a_x - a_y)
    \]
    Otherwise, $r_x + r_y = 51$, so
    \[
    a_x + a_y \equiv r_x + r_y \equiv 51 \equiv 0 \pmod{51}
    \]
    hence
    \[
    51 \mid (a_x + a_y)
    \]

    \medskip

    Hence among every set of $30$ integers, there exist two whose difference or sum is a multiple of $51$.
\end{proof}

\end{document}
